\section{引言}
\label{sec:intro}

将目标检测网络部署到边缘 MPSoC 时，真正决定系统体验的并不只是卷积核的
理论运算量。输入图像需要经过预处理和量化，网络中间张量必须在片上缓冲、
DDR 和处理系统之间移动，不同尺度分支需要保存和重新汇合，最终还要完成
候选框解码与非极大值抑制。即使卷积占据绝大多数乘加，数据布局转换、
partial-PSUM 反馈和层间切换仍可能使一块具备较高峰值吞吐的阵列长期等待。
这一问题在功耗和内存资源受限、但同时集成可编程逻辑与通用处理器的边缘
MPSoC 上尤为突出。传统 GPU/NPU 可以通过大批量和多任务隐藏调度开销，而
batch=1 的实时检测没有这一缓冲空间；每一次空泡都会直接反映为单图延迟。

YOLO 系列以单阶段检测和规则卷积骨干著称\cite{yolov1,yolov2,yolov3}，
YOLOv3-tiny 进一步缩减层数，因此常被用作 FPGA 目标检测加速器的评测网络。
但是，“tiny”并不意味着层形状同质。以固定 416$\times$416 输入为例，完整
双尺度 YOLOv3-tiny 同时包含：空间尺寸从 416$\times$416 到 13$\times$13
的 3$\times$3 卷积；输入通道从 3 到 1024 的变化；1$\times$1 通道压缩；
26$\times$26、$C_{in}=384$ 的分支融合卷积；以及两个 255 通道检测头。
同一固定阵列若按最深层配置，浅层会因空间和通道尾块产生浪费；若按浅层
配置，m13 与 m19 等长 $K$ 层又需要大量 pass 和 partial sum 往返。
\cref{fig:workload} 给出了本文工作负载的 $K$-pass 分布。m0 只需 2 个 pass，
而 m13 和 m19 分别需要 256 和 192 个 pass；固定“加载输入--加载权重--
计算--排出结果”的串行时序无法同时适应两类层。

\begin{figure*}[t]
\centering
\begin{tikzpicture}
\begin{axis}[
  width=0.96\textwidth,height=6.1cm,
  ybar,bar width=7pt,
  ylabel={$K$-pass 数},
  symbolic x coords={m0,m2,m4,m6,m8,m10,m13,m14,m15,m16,m19,P4,P5},
  xtick=data,x tick label style={rotate=35,anchor=east,font=\small},
  ymin=0,ymax=280,grid=major,
  legend style={at={(0.5,1.02)},anchor=south,legend columns=2,draw=none},
  nodes near coords,nodes near coords style={font=\scriptsize,rotate=90,anchor=west},
]
\addplot[fill=LasaBlue!75,draw=LasaBlue] coordinates {
 (m0,2) (m2,8) (m4,16) (m6,32) (m8,64) (m10,128) (m13,256)
 (m14,57) (m15,128) (m16,15) (m19,192) (P4,15) (P5,29)};
\addlegendentry{$K$-pass}
\end{axis}
\node[anchor=north west,align=left,font=\small] at (-0.1,-0.35) {
空间范围：$416^2\rightarrow 13^2$；$C_{out}$：16--1024；\\
检测头 255 通道产生 8 个 Cout 块，最后一块仅 31 lane 有效。};
\end{tikzpicture}
\caption{完整 YOLOv3-tiny 卷积工作负载的异构性。m13 的 256 个 $K$-pass
与 m19 的 192 个 $K$-pass 强调了长层反馈和重放的重要性；早期层则由大空间
tile 与输入搬运主导。}
\label{fig:workload}
\end{figure*}


已有 CNN 加速研究从数据复用、循环变换、片上存储和量化等方向显著提高了
FPGA/ASIC 的卷积效率\cite{kung1982systolic,zhang2015optimizing,
qiu2016going,caffeine2016,chen2016eyeriss,chen2017eyerissv2}。FINN、
hls4ml 等工作进一步展示了定制精度与流式网络映射的可行性
\cite{finn2017,umuroglu2020finn,hls4ml2018}。然而，面向真实目标检测的完整
系统仍存在三类经常被峰值 GOPS 掩盖的损耗。

第一，\textbf{布局与窗口生成损耗}。软件生态和后处理通常采用 HWC 或
NHWC 张量，而阵列计算喜欢 $K\times C_{out}$ 的矩阵片段。若 PS 预先执行
im2col 或逐 pass 重排，DDR 流量和 CPU 时间会随 $K$-pass 数增长；若每个
pass 都重新读取原始 HWC，也会反复搬运同一 tile。FEATHER 强调了数据布局
与数据流协同设计的重要性\cite{feather2024}，本文进一步关注 HWC 输入在
多 pass 阵列中的片上物化和可重放语义。

第二，\textbf{长 $K$ 维反馈损耗}。18 行阵列一次只覆盖 18 个 $K$ 元素，
3$\times$3、512 输入通道的卷积需要 256 个 pass。非末 pass 的 int32 PSUM
必须严格按像素和输出通道顺序写入、再在下一 pass 注入。如果 drain、反馈
读取和下一轮计算完全串行，阵列利用率会受到 pass 数而非 MAC 数限制。

第三，\textbf{上下文切换损耗}。相邻 pass 或相邻 Cout block 之间，权重、
bias、IFM 元数据、PSUM 所有权和输出地址均需切换。简单双缓冲不足以保证
正确性：预载的上下文可能在旧结果尚未提交时覆盖缓存，快速交接也可能在
异常或反压下破坏顺序。因而需要将上下文身份、资源所有权和提交次序纳入
硬件协议，而不是只在性能时间线上画出理想重叠。

这些损耗彼此耦合，使“把某一层做到最快”并不等价于“把完整网络做到最快”。
例如，增大空间 tile 可以减少 dispatch 和边界重启，却同时增加窗口物化深度、
PSUM 容量和长连线压力；把输入预展开为矩阵可以简化阵列接口，却会让软件
重排、DDR 写回和 cache 维护进入关键路径；增加上下文缓冲能够提高预载覆盖，
却必须回答旧上下文何时真正退休、异常反压时谁拥有反馈 bank 等正确性问题。
因此本文围绕三个相互关联的研究问题展开：其一，如何在保持软件可见 HWC
布局的同时，为固定阵列生成可跨 pass 重放的向量；其二，如何让同一组硬件
面对空间型浅层、长 $K$ 深层和分支融合层时选择不同分块，而不修改外部推理
语义；其三，如何在权重、输入与 PSUM 存在真实数据依赖时实现上下文重叠，
并使每一种重叠都能被反压测试、计数器和端到端结果所证伪。

本文所称“完整推理”也有严格边界。它指固定 416 输入下的标准双尺度
YOLOv3-tiny DAG：两个检测头均保留，route 分支、特殊 stride-1 池化、上采样、
重量化拼接和 class-aware NMS 均执行；不把只运行骨干卷积或单一检测头称为
完整网络。与此同时，本文不要求所有算子都进入 PL。MPSoC 的价值在于让
规则、乘加密集的卷积留在可编程逻辑，把低占比但控制不规则的张量操作与
后处理交给 A53。评价对象因此是一个明确的软硬件边界，而不是脱离调度器、
DMA 和数值语义的孤立阵列峰值。

为解决上述问题，本文提出 \textbf{LASA}（Layer-Adaptive HWC-Streaming
Systolic Accelerator）。其核心思想不是让阵列在每一种层上改变物理规模，
而是让 HWC 物化、$K$/Cout/空间分块和上下文调度共同适应层形状：输入 tile
只在片上物化一次，后续 pass 从 URAM/BRAM 缓存重放；32-lane 输出块同时
支持 1$\times$1、3$\times$3 和 255 通道尾块；权重预载、tagged context、
连续 PSUM 和解耦 drain 在保持顺序提交的同时跨越 pass 边界。本文贡献固定
为以下四点。

\begin{enumerate}[leftmargin=2em]
  \item 提出一套\textbf{层自适应脉动阵列映射}。18 路 $K$ 输入与 16 列
  双输出通道 PE 形成 32 通道并行输出，原生覆盖 1$\times$1/3$\times$3、
  任意 $K$ 尾 pass 和非 32 整倍数输出通道。
  \item 提出\textbf{片上 HWC 物化与重放机制}。HWC 输入直接银行化并生成
  卷积窗口，在长层 spatial tile 内只物化一次；后续 $K$-pass 复用片上
  向量，减少 PS 侧重排和重复 DDR 输入。
  \item 提出\textbf{跨 pass 与跨上下文延迟隐藏协议}。tagged context、
  weight preload、fast handoff、continuous PSUM 和静止式 drain ring 将
  权重/反馈/提交与当前计算重叠，并以所有权、不覆盖和顺序提交不变量保证
  反压下的正确性。
  \item 构建\textbf{可验证的完整 MPSoC 推理系统}。完整 13 个卷积在 PL
  运行，A53 处理图结构和非卷积算子；使用逐节点 byte-exact、5000 张
  COCO raw-head 评测、3$\times$1000 性能和 600 s soak 建立从 RTL 到
  系统结果的证据链。
\end{enumerate}

量化不是本文创新。本文使用硬件兼容的 per-tensor INT8 参数固定推理语义，
将 RTL-semantic host 模型作为板端整数结果的权威参照，并分别报告 FP32 和
INT8 COCO 指标。当前部署参数未满足预设精度损失门限，因此模型状态与硬件
部署正确性被明确分离：参数训练可以在不改变 \LASA 硬件的情况下替换，
但不能通过硬件 byte-exact 来掩盖量化精度损失。

为避免实现论文常见的“机制、性能和正确性分别来自不同版本”问题，本文采用
证据优先的方法组织实验。正式 bitstream 必须先通过 200 MHz 的时序、路由、
DRC 和约束覆盖门禁；参数包、输入、软件和输出分别以 SHA256 绑定；板端每个
整数节点先与 RTL-semantic host 比较，再运行 COCO 指标；性能记录必须携带
输出 CRC 与分项计数器。这样，某项优化只有在相同网络语义下输出完全一致，
才有资格进入性能比较。对尚未完成的 CPU 基线和受控 bitstream 消融，正文
保留明确占位并由冻结脚本拒绝投稿状态，而不使用推算值填表。

本文其余部分安排如下：\cref{sec:background} 给出工作负载、循环定义和
设计动机；\cref{sec:system} 描述 MPSoC 系统边界；\cref{sec:microarch}
介绍 \LASA 数据通路；\cref{sec:scheduling} 讨论层自适应分块与跨 pass
调度；\cref{sec:software} 说明完整网络的软件协同；\cref{sec:method}
和 \cref{sec:results} 分别给出实验方法与当前证据；最后讨论限制和结论。
